\begin{lema}[Rearrangement of the \(f6\) contribution]\label{lem:f6} On the grid of side \(s\), let
\(\mathcal I{}\) be the square--line incidence matrix of Section~\ref{sec:rectangulos}, and set
\(\beta_e=\nu_1(e)/s^2\) and \(N_v=\mathcal I{}^{\mathsf T}\operatorname{diag}(v)\mathcal I{}\). Products
of vectors indexed by squares are pointwise. With \(K,L_0,\rho,w\) from Section~\ref{sec:profile-data},
define \begin{align*}
 D&=L_0-K,& V_{ij}&=K_{ij}(W_{ij}-K_{ij}),&
 \Psi&=\rho^{-1}\mathcal I{}^{\mathsf T}(\beta w),\\
 Z&=\Psi+D\rho\Psi,& h&=w-\mathcal I{}Z,& \alpha&=\beta h,\\
 a&=\mathcal I{}^{\mathsf T}\alpha,& b_\beta&=\mathcal I{}^{\mathsf T}\beta,&
 M&=\mathcal I{}D\mathcal I{}^{\mathsf T}.
\end{align*} Here \(W_{ij}\) is the intersection weight, or zero if there is no intersection. Write
\(U(e,j)=\sum_{u\in \familias(e)} K_{uj}-3W(e)\) and \(t(c,e;j)=M_{ce}+\theta_j+U(c,j)+U(e,j)\).
Quadrature of the physical profile, using the generic value at accidental coincidences, gives exactly
\(f6_s=-E_1+E_2\), where \begin{align*}
 E_1={}&\tfrac12\operatorname{tr}(N_\beta D N_\alpha D)
       +\tfrac12 b_\beta^{\mathsf T}Va-3b_\beta^{\mathsf T}D\mathcal I{}^{\mathsf T}(\alpha W)\\
     &-\tfrac12\sum_{f\in D_{\rm obl}}a_f
         [(D N_\beta D)_{ff}+(Vb_\beta)_f],\\
 E_2={}&\sum_c\beta_c\sum_{j\in \familias(c)}\sum_{e\in j}
     \frac{\eta_j\nu_1(e)}s\,t(c,e;j)
     \sum_{u\in \familias(e)\setminus\{j\}}Z_u .
\end{align*} Here \(D_{\rm obl}\) comprises both oblique families. \end{lema} \begin{proof} For fixed
\(c\), the second-order response consists of the local source \(\Lambda\) and the line source \(\omega\).
When we sum the response against \(\beta_e w(e)\), each \(u\in \familias(e)\) receives \([\mathcal
I{}^{\mathsf T}(\beta w)]_u=(\rho\Psi)_u\). By symmetry, the integral part of the response adds \(\rho
D\rho\Psi\). Thus the total mass is \(\rho Z\). The negative source in \(\omega\) integrates \(\Lambda\)
on the fiber of \(u\). On the grid, \(\rho_u\eta_u=1/s\), so the sum subtracts \(\beta_e(\mathcal
I{}Z)_e\) from the weight of \(\Lambda(e;c)\). This leaves mass \(\alpha_e\).

For two generic marks \(c,e\), the mean of \(T_e\) is \(M_{ce}\). Each direct intersection has variance
\(V_{uj}\), since its squared indicator equals the indicator itself. Distinct intersections have distinct
transverse lines in generic position. Thus \(\mathbb E T_e^2=M_{ce}^2+\sum_{u\in \familias(e),j\in
\familias(c)}V_{uj}\). For \(f\in \familias(e)\), the same calculation gives \(\mathbb E
z_f^2=(\sum_{j\in \familias(c)}D_{fj})^2+\sum_{j\in \familias(c)}V_{fj}\). In the sum of
\(\beta_c\alpha_e\mathbb E\Lambda(e;c)\), contracting the indices in the square \(M_{ce}^2\) gives the
trace of \(N_\beta D N_\alpha D\). The variance gives \(b_\beta^{\mathsf T}Va\), and the linear moment
with \(W(e)\) gives \(b_\beta^{\mathsf T}D\mathcal I{}^{\mathsf T}(\alpha W)\). The squares of \(z_f\)
give \((DN_\beta D)_{ff}+(Vb_\beta)_f\), one for each oblique line of \(e\). The coefficients in
\(\Lambda\) give \(E_1\). The matrix \(N_v\) includes the diagonal entries \(\sum_{e\in u}v_e\) even when
the walk returns along the same line.

In the sum over intersections in \(\omega\), fix \(e=u\cap j\), with \(j\in \familias(c)\). The direct
weight times the measure of \(u\) is \(\rho_uK_{uj}=\eta_j\nu_1(e)/s\). Conditioning on availability of
the intersection then marks \(c,e\) on fiber \(j\). In the mean response \(M_{ce}\), we restore the
diagonal term \(\theta_j\). The three intersections with \(j\) at each endpoint now carry the
deterministic weight \(W\) of the marked square. Replacing each original term \(-K\) by \(-W\) adds
\(U(c,j)+U(e,j)\). This leaves \(t(c,e;j)\), multiplied by \(Z_u\) for each transverse line \(u\ne j\) of
\(e\). The intersection source has a negative sign in \(\omega\), and the definition of \(f6\) carries
another negative sign. This gives \(+E_2\) and \(-E_1\). This finite-sum identity leaves the spatial
extrapolation error at status \status{N}. \end{proof}
